\section{Storia dell'informatica}
L'\textbf{informatica}, abbreviazione di informazione automatica è una scienza nata negli anni '60 dalla fusione di vari aspetti scientifici e
ingegneristici. In particolare ha un'origine \textbf{matematica}, dove si è definito che cosa sia "calcolare", e una \textbf{tecnologica}, determinante
gli strumenti utilizzati per il calcolo. Quest'ultimo, nel senso moderno, nasce nell'epoca dell'Illuminismo ed è dovuto, in ambito del Determinismo,
alla necessità di approssimare logaritmi e funzioni trigonometriche mediante polinomi, ottenendo un risultato univoco a partire da date condizioni
iniziali.\par
Formalizzate queste tecniche nel '700, si vide successivamente la necessità di \textbf{meccanizzare} il calcolo. Il primo strumento funzionante fu la
\textbf{macchina analitica}, creata da Charles Babbage e Ada Lovelace nel 19° secolo; proponeva un concetto rivoluzionario: la possibilità di
programmazione per poter ritornare risultati in base ai parametri forniti. Poteva essere usata per sviluppare qualsiasi funzione di complessità e
generalità arbitrarie, ripetendo operazioni sugli stessi elementi.\par
Le origini matematiche dell'informatica ebbero invece luogo con la necessità di verificare se la matematica fosse formalizzabile come insieme finito non
contraddittorio di assiomi. Ciò non è possibile, perché ci si scontra inevitabilmente con le contraddizioni e gli assurdi. L'effettivo limite di ciò
che è calcolabile e non fu definito con il \textbf{teorema di incompletezza} di Gödel; la restrizione del range di studio fu fondamentale per avere
un'idea più chiara del dominio operativo. Andando avanti con gli anni, Alan Turing ideò l'omonima macchina, la cui idea principale era la possibilità
interpretare il \textbf{calcolo} come sequenza anche infinita di sostituzioni di simboli: un algoritmo. Questo era naturalmente un programma finito dove
ogni passo dipende da altrettante informazioni finite. Separa inoltre il controllo dalla memoria, ponendo le basi per quello che sarebbero diventati
i calcolatori moderni. Oltre a ciò, insieme a Church definì la \textbf{Tesi di Church-Turing}, la quale afferma che se un problema è intuitivamente
calcolabile, allora esisterà una macchina di Turing (MDT) o un dispositivo equivalente\footnote{Strumenti come lambda-calcolo, funzioni ricorsive o
linguaggi di programmazione turing completi.} in grado di risolverlo, ovvero calcolarlo. L'ambito di studio dell'informatica è interamente incluso
nell'insieme delle funzioni calcolabili, ovvero i problemi, le cui soluzioni sono fornite dagli algoritmi.

%

\section{Contestualizzazione}
Studiare la teoria della ricorsione e della calcolabilità è utile per fornire consapevolezza in merito ai limiti logici assoluti; insegna in primo luogo
a riconoscere problemi impossibili, fornisce un criterio per capire quando fermarsi e permette di partire da assunzioni più chiare durante lo sviluppo
software.\par
Tutte queste competenze sono in ogni caso usate durante il \textbf{calcolo}, formalmente definito come la risoluzione di un problema tramite un
algoritmo, indipendentemente da quanto tempo o spazio è richiesto per la computazione. Come già menzionato, tuttavia, non sarà sempre possibile trovare
una soluzione per ogni problema proposto, anche se formalmente definibile; per questa ragione, in termini pratici, non esistono programmi che
verifichino la correttezza assoluta o la sicurezza di un dato software. Ciò che più si avvicina allo scenario di controllo ideale è l'\textbf{analisi
del software}, ovvero la verifica del comportamento del programma in base a ciò che ci si aspetta. In particolar modo bisogna garantire: \begin{itemize}
	\item \textbf{Affidabilità}: Il software non deve andare in crash improvvisamente.
	\item \textbf{Sicurezza}: Il software, interfacciandosi con utenti esterni, non deve essere violato o manipolato.
	\item \textbf{Protezione}: Il software, se in controllo di sistemi critici, non deve causare incidenti.
\end{itemize}

\noindent Idealmente, dato un qualunque software con una data proprietà di interesse, esisterebbe un programma di analisi atto al calcolo automatico e
preciso del soddisfacimento della stessa in tempo finito: il debugger ideale. Tuttavia, non è realistico che esista, poiché non abbiamo modo di
verificare a priori se un programma termini o meno; in altre parole non può esistere: \[
	\text{Per ogni programma } p\in \mathbb{L}, \text{ halt}(p) = \text{true} \iff p\text{ termina}
\]

\noindent Ciò vale anche per un eventuale programma antivirus ideale, che segnalerebbe correttamente ogni malware nel caso in cui i relativi programmi
in questione eseguissero compiti malevoli, come rubare o cifrare i dati. Questo perché non è possibile creare un algoritmo universale per un'analisi
infallibile; tenendo a mente la definizione appena discussa, supponiamo un malware $M$ che esegue un secondo programma $P$ su un dato input $x$: se
$P(x)$ termina, un file viene rimosso, altrimenti non arriva mai alla cancellazione. Non potendo sapere se un programma termini a priori, non sarà
altrettanto possibile segnalarlo come malware, mostrando il limite logico.\par
Ancora più difficile sarebbe dimostrare che un software sta agendo in modo malevolo quando il discriminante è una proprietà semantica non banale, come
la cifratura dei dati, nel caso dei ransomware. Da ciò, il \textbf{Teorema di Rice}: non esiste un algoritmo tale che per ogni programma sia capace di
decidere automaticamente una proprietà semantica non banale, come per esempio la terminazione, fornire un risultato specifico, o determinare se due
computazioni siano equivalenti.\par
Ciò non significa che tutto sia precluso; bisogna accettare i limiti imposti dall'insieme dei problemi calcolabili per lavorare esclusivamente su questi
ultimi. In questi termini, non sarà possibile avere analisi automatiche corrette e complete, come anche ottenere risposte sempre precise, data dalla
possibilità del verificarsi di falsi negativi e positivi. In ogni caso, esistono più iter di lavoro per l'analisi del software: \begin{itemize}
	\item \textbf{Model checking}: Controllo esaustivo di un modello del programma da analizzare.
	\item \textbf{Analisi statica}: Controllo automatico ed esaustivo sul programma senza che venga eseguito.
	\item \textbf{Testing}: Verifica di un sottoinsieme delle esecuzioni del sistema.
	\item \textbf{Antivirus firma e sandbox}: Sistema difensivo probabilistico, vulnerabile a stealthing.
\end{itemize}

%

\section{Problema dell'Informatica}
Prima di definire formalmente il problema che l'Informatica si pone di risolvere, è necessario chiarire alcuni concetti preliminari: sappiamo in primo
luogo che il dominio di lavoro è un sottoinsieme dell'insieme dei numeri naturali $\mathbb{N}$ contenente le funzioni calcolabili, le quali sono
definite come un insieme di coppie appartenenti alle parti di $\mathbb{N}$, ovvero: \[
	f\in \mathcal{P}(\mathbb{N}\times \mathbb{N})
\]

\noindent Per conoscere quante sono le funzioni del suddetto dominio, quindi la \textbf{cardinalità} dell'insieme $\mathcal{P}(\mathbb{N}\times
\mathbb{N})$, notiamo che il numero di elementi all'interno di $\mathbb{N}\times\mathbb{N}$ ed $\mathbb{N}$ è equivalente e come conseguenza diretta
sono entrambi numerabili e ordinabili. In questi termini esiste quindi una biezione univoca tra i due insiemi, ovvero che a ogni input corrisponde un
output. La cardinalità dell'insieme delle parti $\mathcal{P}(\mathbb{N}\times \mathbb{N})$ sarà a sua volta equivalente a quella delle funzioni
$\mathbb{N}\times\mathbb{N}$, dimostrando che esistono tante funzione quanti sottoinsiemi di $\mathbb{N}$.\par
Si ritiene ora necessario capire quanti algoritmi risolutivi esistano, definiti in termini logici come sequenze di simboli su un alfabeto $\Sigma$, il
quale deve essere finito e non vuoto; esisteranno quindi tanti programmi quante sono le stringhe componibili con il suddetto alfabeto, il cui insieme
è indicato con $\Sigma^*$. Per esempio: \[
	\Sigma = \{a, b\} \implies \Sigma^* = \{\epsilon, a, b, aa, ab, ba, bb, aaa, ...\}
\]

\noindent Dove $\epsilon$ rappresenta la stringa vuota. Si tratta di un insieme ordinabile privo di un limite superiore e in quanto tale può essere
messo in relazione con quello dei naturali; ne concludiamo che: \[
	|\Sigma^*| = |\mathbb{N}|
\]

\noindent Dunque la cardinalità dell'insieme degli algoritmi è equivalente a quella dell'insieme dei naturali. Inoltre, per il \textbf{teorema di
Cantor}, si può dimostrare come l'insieme delle parti avrà sempre cardinalità maggiore rispetto a quello intero; per esempio: \[
	S = \{a, b\} \implies \mathcal{P}(S) = \{\{a\}, \{b\}, \emptyset, \{a,b\}\} \implies |S| = 2; |\mathcal{P}(S)| = 4
\]

%

\section{Concetti di base richiesti}



\begin{comment}
--- Concetti base
	- f: D\to C		D è dominio, C è codominio
	- f\subseteq D\times C
	- range(f) = {c\in C | \exists d\in D.f(d) = c}, l'insieme degli elementi che possiamo ottenere come output nel codominio, quindi la sua immagine.
	
--- induzione matematica
Sia A \leq \subseteq A\times A relazione antisimmetrica e transitiva, ma non riflessiva.
Per poter avere un'induzione si richiede che \leq sia ben fondata, quindi non esistono catene discendenti infinite (rivedi cosa intende con catena)
Se le catene non sono infinite, significa che ci sono degli elementi minimali, ai quali ci si blocca. Sono i garanti della terminazione. Il loro insieme forma la base di A.

Il principio di induzione dice che:
Dato (A, \leq) ben fondato.
Dato \pi proprietà sugli elementi di A

\forall a\in A, \pi(a) \iff 
	- Passo base: \forall a\in Base(A).\pi(a)
	- Passo induttivo: \forall b < a.\pi(b) \implies \pi(a)
	
Se valgono ambo i passi, l'induzione è dimostrata. In particolare la difficoltà è provare che l'ipotesi induttiva [\forall b < a.\pi(b)] implica la tesi \pi(a)

EG: \forall n\geq 1 \Sigma_{i=1}^n i = \frac{n(n+1)}{2} = \pi(n)
Con A = N, la Base(A) = 0

	- Passo base: n=1
		\pi(1) \implies \Sigma_{i=1}^1 = \frac{1(1+1)}{2} = 1. Corretto.
		
	- Passo induttivo: \forall n \in N {0,1}.\pi(n)
		- Ipotesi induttiva: \forall m < n.\pi(n)
		\sum_{i=1}^n = \sum_{i=1}^{n-1} i+n = \frac{(n+1)(n-1+1)}{2}+n = \frac{n(n-1)}{2}+n = \frac{(n+1)n}{2}

\end{comment}